% ============================================================
\chapter*{Pr\'eface}
\addcontentsline{toc}{chapter}{Pr\'eface}
% ============================================================

\section*{Objectifs du cours}

Ce cours de \textbf{Th\'eorie des Probabilit\'es} s'adresse aux \'etudiants de deuxi\`eme
ann\'ee de licence (L2) en math\'ematiques, informatique ou sciences physiques. Il constitue
une introduction rigoureuse \`a la th\'eorie math\'ematique des probabilit\'es, fond\'ee sur
l'axiomatique de Kolmogorov, tout en maintenant une forte intuition probabiliste.

\medskip

Les probabilit\'es occupent une place centrale dans les math\'ematiques modernes et leurs
applications. De la physique statistique \`a l'apprentissage automatique, de la finance
quantitative \`a la biologie mol\'eculaire, la mod\'elisation probabiliste est devenue un
outil indispensable pour comprendre et quantifier l'incertitude.

\section*{Pr\'erequis}

Pour aborder ce cours dans de bonnes conditions, l'\'etudiant devra ma\^itriser~:
\begin{itemize}
    \item \textbf{Analyse r\'eelle~:} suites, s\'eries, int\'egrales (Riemann), convergence
          uniforme, th\'eor\`emes de convergence domin\'ee (notions de base).
    \item \textbf{Alg\`ebre lin\'eaire~:} espaces vectoriels, matrices, valeurs propres
          (pour le chapitre sur les cha\^ines de Markov).
    \item \textbf{Th\'eorie des ensembles~:} op\'erations ensemblistes, d\'enombrabilit\'e,
          notions de cardinalit\'e.
    \item \textbf{Combinatoire~:} arrangements, combinaisons, principe de multiplication.
\end{itemize}

\section*{Organisation du cours}

Le cours est structur\'e en douze chapitres, organis\'es selon une progression logique~:

\begin{description}
    \item[Chapitres 1--2~:] \textbf{Fondements.} On y d\'efinit l'espace probabilis\'e
    $(\Omega, \mathcal{F}, \PP)$, les axiomes de Kolmogorov, la probabilit\'e conditionnelle
    et la notion d'ind\'ependance. Ces deux chapitres posent le cadre axiomatique sur lequel
    repose tout l'\'edifice.

    \item[Chapitres 3--5~:] \textbf{Variables al\'eatoires.} On introduit les variables
    al\'eatoires discr\`etes puis continues, avec leurs lois, fonctions de r\'epartition,
    densit\'es. Le chapitre~5 d\'eveloppe les notions d'esp\'erance, variance et moments,
    outils fondamentaux du calcul probabiliste.

    \item[Chapitres 6--8~:] \textbf{Lois et transformations.} Le chapitre~6 pr\'esente un
    catalogue des lois classiques (binomiale, Poisson, normale, exponentielle, gamma\ldots).
    Le chapitre~7 \'etudie les fonctions de variables al\'eatoires, et le chapitre~8
    introduit les lois jointes, la covariance et la corr\'elation.

    \item[Chapitres 9--11~:] \textbf{Th\'eor\`emes limites.} C'est le c\oe ur de la
    th\'eorie~: la loi des grands nombres (faible et forte) au chapitre~9, le th\'eor\`eme
    central limite au chapitre~10, et les fonctions caract\'eristiques au chapitre~11,
    outil puissant qui permet de d\'emontrer \'el\'egamment le TCL.

    \item[Chapitre 12~:] \textbf{Introduction aux cha\^ines de Markov.} Ce dernier
    chapitre ouvre la porte aux processus stochastiques en \'etudiant les cha\^ines de
    Markov \`a espace d'\'etats fini ou d\'enombrable.
\end{description}

\section*{M\'ethodologie et notation}

Tout au long de ce cours, nous adoptons les conventions suivantes~:

\begin{itemize}
    \item Les \textbf{d\'efinitions} sont encadr\'ees et num\'erot\'ees. Chaque concept
    nouveau est d'abord d\'efini rigoureusement avant d'\^etre illustr\'e par des exemples.

    \item Les \textbf{th\'eor\`emes} sont \'enonc\'es avec pr\'ecision et, dans la mesure
    du possible, d\'emontr\'es compl\`etement. Lorsqu'une d\'emonstration d\'epasse le
    cadre du cours, nous la rempla\c{c}ons par une esquisse de preuve ou une r\'ef\'erence.

    \item Les \textbf{exemples} sont nombreux et d\'etaill\'es. Ils constituent un
    compl\'ement essentiel \`a la th\'eorie abstraite.

    \item Les \textbf{exercices} en fin de section permettent de v\'erifier la
    compr\'ehension et de d\'evelopper les techniques de calcul.
\end{itemize}

\paragraph{Notations principales.}
\begin{center}
\renewcommand{\arraystretch}{1.3}
\begin{tabular}{cl}
    \toprule
    \textbf{Symbole} & \textbf{Signification} \\
    \midrule
    $\Omega$ & Univers (ensemble des r\'esultats possibles) \\
    $\mathcal{F}$ & Tribu (ensemble des \'ev\'enements) \\
    $\PP$ & Mesure de probabilit\'e \\
    $\E[X]$ & Esp\'erance de la variable al\'eatoire $X$ \\
    $\Var(X)$ & Variance de $X$ \\
    $\Cov(X,Y)$ & Covariance de $X$ et $Y$ \\
    $F_X$ & Fonction de r\'epartition de $X$ \\
    $f_X$ & Densit\'e de $X$ (cas continu) \\
    $\varphi_X(t)$ & Fonction caract\'eristique de $X$ \\
    $X \sim \mathcal{L}$ & $X$ suit la loi $\mathcal{L}$ \\
    $\xrightarrow{\text{p.s.}}$ & Convergence presque s\^ure \\
    $\xrightarrow{\PP}$ & Convergence en probabilit\'e \\
    $\xrightarrow{\mathcal{L}}$ & Convergence en loi \\
    \bottomrule
\end{tabular}
\end{center}

\section*{Perspective historique}

La th\'eorie des probabilit\'es a une histoire riche, qui m\'erite d'\^etre bri\`evement
\'evoqu\'ee. Les premiers travaux remontent \`a la correspondance entre \textbf{Blaise Pascal}
et \textbf{Pierre de Fermat} en 1654, \`a propos du probl\`eme des partis. Au
XVIII\textsuperscript{e} si\`ecle, \textbf{Jakob Bernoulli} d\'emontre la premi\`ere
version de la loi des grands nombres dans son \textit{Ars Conjectandi} (1713), tandis que
\textbf{Abraham de Moivre} obtient une forme primitive du th\'eor\`eme central limite.

\medskip

Au XIX\textsuperscript{e} si\`ecle, \textbf{Pierre-Simon de Laplace} synth\'etise
l'ensemble des connaissances probabilistes de son \'epoque dans sa \textit{Th\'eorie
analytique des probabilit\'es} (1812). \textbf{Pafnouti Tchebychev}, \textbf{Andre\"i Markov}
et \textbf{Aleksandr Liapounov} d\'eveloppent les outils analytiques qui m\`eneront aux
th\'eor\`emes limites modernes.

\medskip

La r\'evolution d\'ecisive survient en 1933, quand \textbf{Andre\"i Kolmogorov} publie ses
\textit{Grundbegriffe der Wahrscheinlichkeitsrechnung}, o\`u il fonde la th\'eorie des
probabilit\'es sur la th\'eorie de la mesure. Cette axiomatisation, que nous adoptons dans
ce cours, a permis de r\'esoudre de nombreux paradoxes et a ouvert la voie \`a la th\'eorie
moderne des processus stochastiques.

\section*{Conseils aux \'etudiants}

\begin{enumerate}
    \item \textbf{Travaillez r\'eguli\`erement.} Les probabilit\'es ne s'apprennent pas
    la veille de l'examen. Chaque chapitre s'appuie sur les pr\'ec\'edents.

    \item \textbf{Faites les exercices.} La compr\'ehension passive de la th\'eorie ne
    suffit pas. C'est en r\'esolvant des probl\`emes que l'on d\'eveloppe l'intuition
    probabiliste.

    \item \textbf{Dessinez.} Repr\'esentez les ensembles, tracez les densit\'es, esquissez
    les fonctions de r\'epartition. La visualisation est un alli\'e pr\'ecieux.

    \item \textbf{V\'erifiez les cas particuliers.} Lorsque vous obtenez une formule
    g\'en\'erale, testez-la sur des exemples simples. Cela permet de d\'etecter les erreurs
    et de renforcer la compr\'ehension.

    \item \textbf{Distinguez le discret du continu.} De nombreuses erreurs proviennent
    de la confusion entre sommes et int\'egrales, entre loi de masse et densit\'e. Soyez
    toujours conscient du cadre dans lequel vous travaillez.
\end{enumerate}

\section*{R\'ef\'erences principales}

Ce cours s'inspire de nombreux ouvrages classiques, parmi lesquels~:
\begin{itemize}
    \item J.~Jacod et P.~Protter, \textit{Probability Essentials}, Springer.
    \item P.~Billingsley, \textit{Probability and Measure}, Wiley.
    \item W.~Feller, \textit{An Introduction to Probability Theory and Its Applications},
          volumes~I et~II, Wiley.
    \item J.-F.~Le Gall, \textit{Int\'egration, Probabilit\'es et Processus Al\'eatoires},
          \'Ecole Normale Sup\'erieure.
    \item D.~Foata et A.~Fuchs, \textit{Calcul des Probabilit\'es}, Dunod.
    \item A.~Raugi, \textit{Cours de probabilit\'es}, Universit\'e de Rennes.
    \item G.~Grimmett et D.~Stirzaker, \textit{Probability and Random Processes}, Oxford.
\end{itemize}

\bigskip

\hfill\textit{Bonne lecture et bon travail \`a tous.}
